\chapter{数据集自主选取}

\section{课程数据集要求说明}

本章对应课程评分标准中的“数据集自主选取”。根据课程考查要求，最终报告数据集需要满足以下条件：第一，数据来源应来自 Kaggle、Hugging Face、阿里天池、和鲸或 OpenML 等公开平台；第二，发布时间或更新年份应为 2024--2025 年；第三，数据格式应为 CSV 结构化表格数据；第四，样本量不少于 1000 行；第五，有效特征不少于 8 个；第六，数据应天然支持回归、分类、聚类中至少两类机器学习任务。

这些要求的目的，是避免直接复用课堂小样例或来源不明数据，使课程报告能够体现自主选题、自主数据筛选和完整建模流程。本文后续所有实验均围绕最终选定的数据集展开，不再使用旧实验中的数据作为主数据集。

\section{候选数据集比较与 pfm 数据集弃用原因}

在候选数据比较阶段，曾考虑过 \texttt{pfm\_train.csv} 和 \texttt{pfm\_test.csv}。该数据集适合员工离职预测等分类任务，但不作为本报告主数据集，主要原因包括：来源平台和发布/更新时间不如当前数据集稳定；主题与已有实验三员工离职预测高度相似，容易导致最终报告与旧实验报告雷同；该数据主要天然支持分类任务，回归任务设计不够自然；从课程期末考查角度看，继续使用该类数据难以体现新的分析视角。

因此，本文最终选择与数字生活方式相关的 2025 Digital Lifestyle Benchmark Dataset。该数据集围绕设备使用、社交媒体、通知、睡眠、学习、活动、风险标记、生产力得分和数字依赖得分展开，既能支持分类和回归，也适合进行聚类画像分析。

\section{最终数据集来源与合规性检查}

最终数据集为 Digital Lifestyle Benchmark Dataset (2025)，来源于 Hugging Face。表 \ref{tab:dataset-compliance} 按课程要求逐项列出合规性检查结果。可以看到，该数据集在来源平台、年份、CSV 格式、样本量、字段数量和任务支持方面均满足课程要求。

\begin{table}[H]
\centering
\caption{数据集合规性与基本信息汇总}
\label{tab:dataset-compliance-summary}
\begin{tabular}{ll}
\toprule
项目 & 内容 \\
\midrule
数据集名称 & 2025 Digital Lifestyle Benchmark Dataset \\
来源页面 & Hugging Face 数据集页面 \\
数据格式 & CSV \\
样本量 & 3500 条记录 \\
字段数 & 24 个字段 \\
许可证 & CC BY 4.0 \\
随机种子 & 42 \\
运行约束 & CPU only，非深度学习 \\
\bottomrule
\end{tabular}
\end{table}


需要说明的是，该数据集是合成教学/benchmark 数据。本文将其用于课程建模训练和行为数据分析，不将其解释为真实人群医学诊断数据，也不据此给出个体健康判断。

\section{数据字段、样本规模与任务支持情况}

原始数据包含 3500 条记录和 24 个字段。字段大致分为四类：人口背景字段，例如年龄、性别、地区、收入、教育、日常角色和设备类型；数字行为字段，例如每日设备时长、手机解锁次数、通知数量和社交媒体使用时间；生活习惯字段，例如学习时间、身体活动天数、睡眠时长和睡眠质量；结果字段，例如 \texttt{high\_risk\_flag}、\texttt{productivity\_score} 和 \texttt{digital\_dependence\_score}。

基于这些字段，本文设置三类任务：分类任务预测 \texttt{high\_risk\_flag}，回归任务预测 \texttt{digital\_dependence\_score} 并检验 \texttt{productivity\_score}，聚类任务根据数字行为和生活习惯形成用户画像。与 pfm 数据集相比，该数据集更贴合期末报告的综合建模要求，也更适合体现分类、回归、聚类三类机器学习任务。
